\documentclass[12pt,a4paper]{article}
\input{../../.common/preamble.sty}

\title{\textbf{Han cantado línea}}

\begin{document}
\maketitle

Dado un \textbf{fichero de texto} y un \textbf{número de línea}, encuentra la línea que corresponde.

\subsection*{Notas}

\begin{itemize}
    \item Si el número de línea indicado no corresponde con ninguno del fichero se debe devolver \mintinline{python}{None}.
    \item Ten en cuenta que la primera línea debe empezar en 1.
\end{itemize}

\subsection*{Ejemplo}

Supongamos que el fichero contiene el siguiente texto:

\begin{minted}[
    linenos,
    numbersep=5pt,
    frame=single,
    framesep=3mm,
]{text}
Beautiful is better than ugly
Explicit is better than implicit
Simple is better than complex
Complex is better than complicated
Flat is better than nested
Sparse is better than dense
Readability counts
Special cases arent special enough to break the rules
Although practicality beats purity
Errors should never pass silently
Unless explicitly silenced
In the face of ambiguity, refuse the temptation to guess
There should be one—and preferably only one—obvious way to do it
Although that way may not be obvious at first unless you're Dutch
Now is better than never
Although never is often better than right now
If the implementation is hard to explain, it's a bad idea
If the implementation is easy to explain, it may be a good idea
Namespaces are one honking great idea—let's do more of those
\end{minted}

Si el número de línea es 3, la salida sería:

\begin{minted}{text}
Simple is better than complex
\end{minted}

\end{document}
